%! TEX root = ../am2-18-19.tex

\chapter{Ecuații diferențiale de ordin superior}

\section{Ecuații rezolvabile prin cuadraturi}

\indent\indent Acestea sînt ecuații care se pot rezolva prin integrări
succesive. Astfel, forma lor generală este $ y^{(n)} = f(x) $, în varianta
neomogenă. Soluția generală va depinde de $ n $ constante arbitrare,
rezultate în urma integrărilor succesive.

\textbf{Exemplu:}
\[
  y^{\prime\prime} = \arcsin x + \frac{x}{\sqrt{1 - x^2}} - \frac{\pi}{6}.
\]
\emph{Soluție:} Integrăm succesiv și obținem:
\begin{align*}
  y^{\prime} &= x \arcsin x - \dfrac{\pi}{6} x + c_1 \\
  y &= \dfrac{x^2}{2} \arcsin x + \dfrac{1}{4}x \sqrt{1 - x^2} - %
  \dfrac{1}{4} \arcsin x - \dfrac{\pi}{12} x^2 + c_1 x + c_2.
\end{align*}

Dacă nu se dau condiții inițiale, soluția rămîne exprimată ca mai sus,
adică în funcție de constantele $ c_1 $ și $ c_2 $.

\vspace{1cm}

Cazul omogen se rezolvă și mai simplu: dacă $ y^{(n)} = 0 $, atunci
$ y $ este un polinom de $ x $ de gradul $ n - 1 $.

%%%%%%%%%%%%%%%%%%%%%%%%%%%%%%%%%%%%%%%%%%%%%%%%%%%%%%%%%%%%%%%%%%%%%%

\section{Ecuații de forma $ f(x, y^{(n)}) = 0 $}

Avem două cazuri care trebuie tratate distinct:
\begin{enumerate}[(a)]
  \item Dacă ecuația poate fi rezolvată în raport cu $ y^{(n)} $, adică dacă
    $ \dfrac{\partial f}{\partial y^{(n)}} \neq 0 $, atunci obținem una
      sau mai multe ecuații ca în secțiunea anterioară;
    \item Dacă ecuația nu este rezolvabilă cu ajutorul funcțiilor elementare
      în raport cu $ y^{(n)} $, dar cunoaștem o reprezentare parametrică
      a curbei $ f(u, v) = 0 $, adică $ u = g(t), v = h(t) $, cu
      $ t \in [\alpha, \beta] $, atunci soluția generală se poate da sub formă
      parametrică:
      \[
        \begin{cases}
          x &= g(t) \\
          dy^{(n-1)} &= y^{(n)} dx = h(t) g^{\prime}(t) dt
        \end{cases},
      \]
      de unde putem obține succesiv:
      \[
        \begin{cases}
          y^{(n-1)}  &= h_1(t, c_1) \\
             \dotfill \\
          y(t) &= h_n(t, c_1, \dots, c_n).
        \end{cases}
      \]
\end{enumerate}

\textbf{Exemplu:}
\[
  x - e^{y^{\prime\prime}} + y^{\prime\prime} = 0.
\]
\emph{Soluție:} Putem nota $ y^{\prime\prime} = t $ și atunci $ x(t) = e^t - t $.
Deoarece avem $ dy^{\prime} = y^{\prime\prime} dx $, rezultă:
\[
  dy^{\prime} = t(e^t - 1) dt \Rightarrow y^{\prime} = -\dfrac{t^2}{2} + te^t - e^t + c_1.
\]
Mai departe, $ dy = y^{\prime} dx $, deci:
\[
  dy = \Big( - \dfrac{t^2}{2} + te^t - e^t + c_1 \Big)(e^t - 1) dt.
\]
În fine, soluția este:
\[
  y(t) = e^{2t} \Big( \dfrac{t}{2} - \dfrac{3}{4} \Big) + e^t \Big( - \dfrac{t^2}{2} + %
  1 + c_1 \Big) + \dfrac{t^3}{6} - c_1 t + c_2.
\]

%%%%%%%%%%%%%%%%%%%%%%%%%%%%%%%%%%%%%%%%%%%%%%%%%%%%%%%%%%%%%%%%%%%%%%

\section{Ecuații de forma $ f(y^{(n-1)}, y^{(n)}) = 0 $}

Din nou, distingem două cazuri:
\begin{enumerate}[(a)]
  \item Dacă ecuația este rezolvabilă prin funcții elementare în raport cu $ y^{(n)} $,
    atunci putem nota $ z = y^{(n-1)} $ și avem $ z^{\prime} = f(z) $. Ecuația devine cu
    variabile separabile pentru $ z $ și se rezolvă corespunzător, conducînd la
  $ z = y^{(n-1)} = f_1(x, c_1) $, care este de tipul anterior;
  \item Dacă ecuația nu este rezolvabilă prin funcții elementare în raport cu
    $ y^{(n)} $, dar cunoaștem o reprezentare parametrică de forma:
    \[
      y^{(n-1)} = h(t), \quad y^{(n)} = g(t), \quad t \in [\alpha, \beta],
    \]
    atunci putem folosi $ dy^{(n-1)} = y^{(n)} dx $, pentru a obține soluția
    pas cu pas, sub formă parametrică:
    \begin{align*}
      x &= \ds\int\frac{h^{\prime}}{g} dt + c_1 \\
      y^{(n-1)} &= h(t) \\
      dy^{(n-2)} &= h(t) = \dfrac{hh^{\prime}}{g} dt \\
      y^{(n-2)} &= \ds\int \frac{hh^{\prime}}{g} dt + c_2 \\
                &\vdots \\
      y &= \int y^{\prime} dx + c_n.
    \end{align*}

\end{enumerate}

\textbf{Exemplu:} $ y^{\prime\prime} = - \sqrt{1 - y^{^{\prime}2}} $.
\emph{Soluție:} Facem schimbarea de funcție $ z(x) = y^{\prime} $ și ecuația devine:
\[
  - \dfrac{dz}{\sqrt{1 - z^2}} = dx, \quad |z| < 1.
\]
Soluția generală este: $ \arccos z = x + c_1 $, deci $ z = \cos(x + c_1) $.

Mai departe, obținem $ y(x) = \sin(x + c_1) + c_2 $.

De asemenea, trebuie să mai remarcăm și soluțiile particulare $ y = \pm x + c $.

%%%%%%%%%%%%%%%%%%%%%%%%%%%%%%%%%%%%%%%%%%%%%%%%%%%%%%%%%%%%

\section{Ecuații de forma $ f(x, y^{(k)}, \dots, y^{(n)}) = 0 $}

Ecuația se rezolvă cu schimbarea de funcție $ y^{(k)} = z(x) $ și rezultă
o ecuație de ordin $ n - k $:
\[
  f(x, z, z^{\prime}, \dots, z^{(n-k)}) = 0,
\]
pe care o integrăm.

\textbf{Exemplu:} $ (1 + x^2)y^{\prime\prime} = 2xy^{\prime} $.

\emph{Soluție:} Cu schimbarea de funcție $ y^{\prime} = z(x) $, obținem ecuația:
\[
  \dfrac{z}{z^{\prime}} = \dfrac{2x}{1 + x^2},
\]
iar apoi, prin integrare, rezultă $ z = y^{\prime} = c_1(1 + x^2) $. În fine:
\[
  y(x) = c_1 x + c_1 \dfrac{x^3}{3} + c_2.
\]
\vspace{1cm}

\textbf{Exemplu 2:} $ y \cdot y^{\prime\prime} - yy^{^{\prime}2} = 0 $.

\emph{Soluție:} Notăm $ y^{\prime} = z $ și obținem:
\[
  y\cdot z^{\prime} - yz^2 = 0 \Rightarrow y(z^{\prime} - z^2) = 0,
\]
de unde rezultă $ y = 0 $ sau $ z^{\prime} - z^2 = 0 $.

Din a doua variantă, deducem succesiv:
\[
  \dfrac{dz}{dx} = z^2 \Rightarrow \dfrac{dz}{z^2} = dx \Rightarrow %
  -\dfrac{1}{z} = x + c_1.
\]
Mai departe:
\[
  -\dfrac{1}{y^{\prime}} = x + c_1 \Rightarrow y^{\prime} = \dfrac{-1}{x + c_1} %
  \Rightarrow y = -\ln(x + c_1) + c_2.
\]


%%%%%%%%%%%%%%%%%%%%%%%%%%%%%%%%%%%%%%%%%%%%%%%%%%%%%%%%%%%%%%%%%%%%%%

\section{Ecuații de forma $ f(y^{\prime}, y^{\prime\prime}, \dots, y^{(n)}) = 0 $}

Să remarcăm că astfel de ecuații nu conțin pe $ y $. Deci putem lua pe $ y^{\prime} $
ca variabilă independentă și pe $ y^{\prime\prime} $ ca fiind funcție de $ y^{\prime} $.
Astfel, reducem discuția la un caz anterior.

\textbf{Exemplu:} $ x^2 y^{\prime\prime} = y^{^{\prime}2} - 2xy^{\prime} + 2x^2 $.

\emph{Soluție:} Facem schimbarea de funcție $ y^{\prime} = z(x) $ și obținem o
ecuație Riccati:
\[
  z^{\prime} = \dfrac{z^2}{x^2} - 2 \cdot \dfrac{z}{x} + 2.
\]
Observăm soluția particulară $ z_p = x $ și integrăm, cu $ z = x + \dfrac{1}{u(x)} $.

Obținem succesiv:
\begin{align*}
  z(x) &= x + \dfrac{c_1x}{x + c_1} \Rightarrow \\
  y^{\prime}(x) &= x + \dfrac{c_1x}{x + c_1} \Rightarrow \\
  y(x) &= \dfrac{x^2}{2} + c_1 x - c_1^2 \ln | x + c_1 | + c_2,
\end{align*}

%%%%%%%%%%%%%%%%%%%%%%%%%%%%%%%%%%%%%%%%%%%%%%%%%%%%%%%%%%%%%%%%%%%%%%

\section{Ecuații autonome (ce nu conțin pe $ x $)}

În cazul acestor ecuații, putem micșora ordinul cu o unitate, notînd
$ y^{\prime} = p $ și luăm pe $ y $ variabilă independentă.

\begin{remark}\label{rk:aut-const}
  Există posibilitatea de a pierde soluții de forma $ y = c $ prin această
  metodă, deci trebuie verificat dacă este cazul separat.
\end{remark}

\textbf{Exemplu:} $ 1 + y^{^{\prime}2} = 2yy^{\prime} $.
\emph{Soluție:} Putem lua $ y^{\prime} = p $ drept funcție, iar pe $ y $ drept
variabilă independentă. Rezultă:
\[
  y^{\prime\prime} = \dfrac{d}{dx} \Big( \dfrac{dy}{dx} \Big) = p \dfrac{dp}{dy}.
\]
Atunci ecuația devine:
\[
  \dfrac{2p dp}{1 + p^2} = \dfrac{dy}{y} \Rightarrow %
  y = c_1(1 + p^2).
\]
Acum trebuie să obținem pe $ x $ ca funcție de $ p $ și $ c_1 $.
Deoarece $ dx = \dfrac{1}{p} dy $, iar $ dy = 2c_1 pdp $, rezultă
$ dx = 2c_1 dp $. Deci $ x = 2c_1p + c_2 $, iar soluția generală devine
$ x(p) = 2c_1 p + c_2 $. Din expresia lui $ y $ de mai sus, rezultă:
\[
  y = c_1 + \dfrac{(x - c_2)^2}{4c_1}.
\]

\vspace{1cm}

\textbf{Exemplu 2:} $ y^{\prime\prime} = y^{^{\prime}2} = 2e^{-y} $.

\emph{Soluție:} Notăm $ y^{\prime} = p $ și luăm ca necunoscută $ p = p(y) $.
Rezultă:
\[
  y^{\prime\prime} = \dfrac{dy^{\prime}}{dx} = p^{\prime}p \Rightarrow p^{\prime}p + p^2 = 2e^{-y}.
\]
Dacă notăm $ p^2 = z $, obținem o ecuație liniară neomogenă:
\[
  z^{\prime} + 2z = 4e^{-y},
\]
ce are ca soluție generală $ z(y) = c_1 e^{-2y} + 4e^{-y} $.

Revenim la $ y $ și avem:
\[
  z = p^2 = y^{^{\prime}2} \Rightarrow y^{\prime} = \pm \sqrt{c_1 e^{-2y} + 4e^{-y}},
\]
adică:
\[
  \dfrac{dy}{\pm \sqrt{c_1 e^{-2y} + 4e^{-y}}} = dx \Rightarrow %
  x + c_2 = \pm \dfrac{1}{2} \sqrt{c_1 + 4e^y}.
\]
Echivalent, putem scrie soluția și în forma implicită:
\[
  e^y + \dfrac{c_1}{4} = (x + c_2)^2.
\]

%%%%%%%%%%%%%%%%%%%%%%%%%%%%%%%%%%%%%%%%%%%%%%%%%%%%%%%%%%%%%%%%%%%%%%

\section{Ecuații Euler}

O ecuație Euler are forma generală:
\[
  \sum_{k = 0}^n a_k x^k y^{(k)} = f(x),
\]
pentru varianta omogenă și $ f(x) = 0 $ pentru varianta neomogenă, cu $ a_i $
constante reale.

Ecuațiile Euler se pot transforma în ecuații cu coeficienți constanți prin notația
(schimbarea de variabilă) $ |x| = e^t $.

De remarcat faptul că, deoarece funcția necunoscută este $ y = y(x) $, pentru
a obține $ y^{\prime} = \dfrac{dy}{dx} $ trebuie să aplicăm regula de derivare
a funcțiilor compuse. Folosind notația de la fizică $ \dot{y} = \dfrac{dy}{dt} $,
putem scrie:
\[
  \dfrac{dy}{dx} = \dfrac{dy}{dt} \cdot \dfrac{dt}{dx} = \dot{y} \cdot e^{-t}.
\]
Mai departe, de exemplu, pentru derivatele superioare:
\begin{align*}
  \dfrac{d^2 y}{dx^2} &= \dfrac{d}{dt} \Big( \dot{y} \cdot e^{-t} \Big) %
  \cdot \dfrac{dt}{dx} \\
                      &= e^{-2t} (\ddot{y} - \dot{y}).
\end{align*}

\vspace{1cm}

\textbf{Exemplu:} $ x^2 y^{\prime\prime} + xy^{\prime} + x = x $.

\emph{Soluție:} Facem substituția $ |x| = e^t $ și, folosind calculele
de mai sus, obținem:
\[
  e^{2t} \cdot e^{-2t} (\ddot{y} - \dot{y}) - e^t \cdot e^{-t} + y(t) = e^t.
\]
Echivalent: $ \ddot{y} - 2\dot{y} + y = e^t $.

Aceasta este o ecuație liniară de ordinul al doilea, neomogenă. Asociem ecuația
algebrică $ r^2 - 2r + 1 = (r - 1)^2 = 0 $, deci soluția generală a variantei
omogene este:
\[
  y(t) = c_1 e^t + c_2 te^t.
\]

Folosind metoda variației constantelor, obținem succesiv:
\begin{align*}
  y(t) &= c_1 e^t + c_2 t e^t + \dfrac{t^2 e^t}{2} \Rightarrow \\
  y(x) &= c_1 x + c_2 x \ln x + \dfrac{x \ln^2 x}{2}.
\end{align*}


%%%%%%%%%%%%%%%%%%%%%%%%%%%%%%%%%%%%%%%%%%%%%%%%%%%%%%%%%%%%%%%%%%%%%%

\section{Exerciții}

1. Rezolvați următoarele ecuații Euler:
\begin{enumerate}[(a)]
  \item $ x^2 y^{(2)} - 3xy^{\prime} + 4y = 5x, \quad x > 0 $;
  \item $ x^2 y^{(2)} - xy^{\prime} + y = x + \ln x, \quad x > 0 $;
  \item $ 4(x + 1)^2 y^{(2)} + y = 0, x > -1 $;
  \item $ xy^{\prime\prime} + 3y^{\prime} = 1 $;
  \item $ x^3 y^{\prime\prime\prime} + x^2y^{\prime\prime} - 2xy^{\prime} + 2y = \dfrac{1}{x^2} $.
\end{enumerate}


\vspace{1cm}

2. Rezolvați următoarele ecuații:
\begin{enumerate}[(a)]
  \item $ y^{(3)} = \sin x + \cos x $;
  \item $ y^{(3)} + y^{(2)} = \sin x $;
  \item $ xy^{(2)} + (x - 2)y^{(1)} - 2y = 0 $;
  \item $ x^2 y^{(2)} = y^{\prime 2} $;
  \item $ y^{(4)} - 3y^{(3)} + 3y^{(2)} - y^{(1)} = 0 $;
  \item $ y^{(2)} - y^{(1)} - 2y = 3e^{2x} $;
  \item $ 4(x + 1)^2 y^{(2)} + y = 0 $;
  \item $ \begin{cases}
      x^\prime &= x + ye^{2t} \\
      y^\prime &= y - xe^{2t} + 1
    \end{cases} $,
    cu condițiile inițiale $ x(0) = 1 $ și $ y(0) = 0 $.
\end{enumerate}

Amintim, pentru ecuații liniare și neomogene de ordinul întîi, cu forma generală:
\[
  y' = P(x) y + Q(x),
\]
avem soluțiile:
\begin{align*}
  y_p &= c \cdot \exp\Big( - \int_{x_0}^x P(t) dt \Big) \\
  y_g &= \Big( c + \int Q(t) \cdot \exp \Big( \int_{x_0}^t P(s) ds \Big) dt \Big) \cdot \exp \Big(- \int_{x_0}^x P(t) dt \Big).
\end{align*}

\emph{Indicații:}
\begin{enumerate}[(a)]
  \item Integrări succesive;
  \item Notăm $ y^{(2)} = z $, deci ecuația devine $ z^\prime + z = \sin x $, care este
    o ecuație liniară și neomogenă, de ordinul întîi. Soluția generală este:
    \[
      z(x) = e^{-x}\Big( c + \dfrac{e^x}{2} (\sin x + \cos x) \Big),
    \]
    iar apoi $ y^\prime(x) = \ds\int z(x) dx $ și $ y(x) = \ds\int y^\prime (x) dx $ etc.
  \item Notăm $ z = y^\prime + y $, deci ecuația devine $ xz^\prime - 2z = 0 $, care
    are soluția evidentă $ z = c_1 x^2 $, din care obtinem $ y^\prime + y = c_1 x^2 $, care
    este o ecuație liniară și neomogenă, cu soluția generală:
    \[
      y = c_2 e^{-x} + c_1 x^2 - 2c_1 x + 2c_1.
    \]
  \item Notăm $ y^\prime = z $ și obținem $ x^2 z^\prime = z^2 $. Observăm soluția particulară
    $ z = 0 $, deci $ y = c $, iar în celelalte cazuri, avem:
    \[
      x^2 \dfrac{dz}{dx} = z^2 \Rightarrow \dfrac{dx}{x^2} = \dfrac{dz}{z^2}.
    \]
    Aceasta conduce la $ \dfrac{1}{z} = \dfrac{1}{x} + c $ și ne întoarcem
    la $ y $:
    \[
      dx = dy \Big( \dfrac{1}{x} + c \Big) \Rightarrow \dfrac{xdx}{cx + 1} = dy,
    \]
    care, prin integrare, conduce la:
    \[
      y = \dfrac{x}{c} - \dfrac{1}{c^2} \ln(cx + 1) + c_1,
    \]
    pentru $ c \neq 0 $ și $ y^\prime = x $ dacă $ c = 0 $, adică $ y = \dfrac{x^2}{2} + c_1 $.
  \item Ecuația caracteristică este:
    \[
      r^4 - 3r^3 + 3r^2 - r = 0 \Rightarrow r(r^3 - 3r^2 + 3r - 1) = 0,
    \]
    de unde observăm soluțiile $ r = 0, r = 1 $ etc.
  \item Ecuația caracteristică este $ r^2 - r - 2 = (r + 1)(r - 2) = 0 $.
  \item Este o ecuație Euler în raport cu $ t = x + 1 $.
  \item Aplicăm metoda substituției (eliminării).
\end{enumerate}

%%% Local Variables:
%%% mode: latex
%%% TeX-master: "../am2-18-19"
%%% End:
